% Semantic contract: four phases, left to right, each closing exactly one
% attack surface the previous phase left open. The arrow between phase N and
% N+1 carries the reason N alone was insufficient ("needs:"); the box under
% each phase names what that phase forecloses; the bottom strip is what has
% actually been machine-checked, not what the design merely claims. Absolute
% coordinates throughout (rather than relative positioning) so that every
% row's vertical budget is chosen deliberately, not inferred from whichever
% cell happens to wrap the most.
\begin{figure}[H]
\centering
\resizebox{0.98\linewidth}{!}{%
\begin{tikzpicture}[
  pd figure,
  phase/.style={draw=hhink,line width=.62pt,rounded corners=2.5pt,
    inner xsep=6pt,inner ysep=6pt,align=left,text width=3.05cm,
    font=\footnotesize,minimum height=4.3cm,anchor=north},
  phase13/.style={phase,fill=hhsand!45},
  phase4/.style={phase,fill=hhcobalt!12,draw=hhcobalt!65},
  needs/.style={pd focus fill,rounded corners=5pt,font=\footnotesize,
    inner sep=2.5pt,align=center,text width=2.5cm,anchor=south},
  forecloses/.style={draw=hhcobalt!55,line width=.5pt,fill=hhcobalt!8,
    rounded corners=2pt,inner xsep=5pt,inner ysep=4pt,align=left,
    text width=3.05cm,font=\footnotesize,text=hhink,anchor=north},
  status/.style={pd direct label,align=left,text width=3.05cm,anchor=north,
    fill=none,inner sep=0pt},
]

% ---------- x-centers for the four columns and the three gaps ----------
\def\xa{0}\def\xb{4.05}\def\xc{8.10}\def\xd{12.15}

% ---------- the four phase boxes (tops all at y=0, uniform height) ----------
\node[phase13] (p1) at (\xa,0)
  {\textbf{Phase 1}\\ HS256 Symmetric Pinning\\[3pt]
   Verifier pins the algorithm from policy, not the token header: one
   shared HMAC key per daemon--harbor pair.};
\node[phase13] (p2) at (\xb,0)
  {\textbf{Phase 2}\\ Ed25519 Asymmetric Identity\\[3pt]
   Daemon signs as root CA with a private key; every verifier holds only
   the matching public key.};
\node[phase13] (p3) at (\xc,0)
  {\textbf{Phase 3}\\ Macaroon Stigmergic Delegation\\[3pt]
   Each hop appends a narrower capability set, signs the whole chain
   offline; the verifier checks a monotone subset per hop.};
\node[phase4] (p4) at (\xd,0)
  {\textbf{Phase 4}\\ Cuckoo Revocation\\[3pt]
   An append-only log of revoked \texttt{kid}s is gossiped between
   daemons and projected into a local cuckoo filter for $O(1)$ lookups.};

% ---------- needs: pills, clear of the boxes above and the arrows below ----------
\node[needs] (n12) at ($(\xa,0)!0.5!(\xb,0)+(0,2.05)$)
  {needs: decentralized A2A with no shared HMAC secret};
\node[needs] (n23) at ($(\xb,0)!0.5!(\xc,0)+(0,2.05)$)
  {needs: multi-hop delegation without round-tripping the Daemon};
\node[needs] (n34) at ($(\xc,0)!0.5!(\xd,0)+(0,2.05)$)
  {needs: withdrawing a leaked or policy-reversed card before its TTL};

% ---------- arrows, at box mid-height, with a leader up to each pill ----------
\draw[pd arrow] (p1.east) -- (p2.west);
\draw[pd hairline] ($(\xa,0)!0.5!(\xb,0)$) -- (n12.south);
\draw[pd arrow] (p2.east) -- (p3.west);
\draw[pd hairline] ($(\xb,0)!0.5!(\xc,0)$) -- (n23.south);
\draw[pd arrow] (p3.east) -- (p4.west);
\draw[pd hairline] ($(\xc,0)!0.5!(\xd,0)$) -- (n34.south);

% ---------- "Forecloses:" boxes, tops all at y=-5.0 ----------
\node[forecloses] (f1) at (\xa,-5.0)
  {\textbf{Forecloses:} header-selected algorithm confusion
   (CVE-2015-9235, CVE-2023-48238).};
\node[forecloses] (f2) at (\xb,-5.0)
  {\textbf{Forecloses:} shared-secret theft --- no verifier holds a
   secret that could mint a card.};
\node[forecloses] (f3) at (\xc,-5.0)
  {\textbf{Forecloses:} downstream privilege growth --- each hop's
   subset check is monotone, so no hop widens scope.};
\node[forecloses] (f4) at (\xd,-5.0)
  {\textbf{Forecloses:} post-issue authority --- a leaked or
   disciplined card that should die before its natural TTL.};

% ---------- mechanized-status strip, hairline well below the tallest box ----------
\draw[pd hairline] (\xa,-7.7) -- (\xd,-7.7);
\node[pd panel title,anchor=west] at (\xa,-8.05)
  {Mechanized status (ProVerif 2.05, \S\ref{sec:correct}):};

\node[status,text=hhteal] (s1) at (\xa,-8.45)
  {\textbf{TRUE} --- \texttt{v1\_refined.pv}: secrecy + authentication};
\node[status,text=hhteal] (s2) at (\xb,-8.45)
  {\textbf{TRUE} --- \texttt{v2\_asymmetric.pv}: secrecy + authentication};
\node[status,text=hhteal] (s3) at (\xc,-8.45)
  {\textbf{TRUE} --- \texttt{v3\_delegation.pv}: authentication
   correspondence (attenuation modeled separately, App.~\ref{app:phase3})};
\node[status,text=hhamber] (s4) at (\xd,-8.45)
  {\textbf{not modeled} --- soundness argued by inspection only
   (\S\ref{sec:revocation})};

\end{tikzpicture}%
}
\caption{The four Harbor Card phases, left to right. Each arrow's
\emph{needs:} pill names the reason the phase to its left was not enough,
drawn from this chapter's own account of why the next phase was added; each
\emph{Forecloses:} box names the attack surface that phase closes (Table
\ref{tab:anchor-four-phases} tabulates the same four rows for lookup). The
bottom strip is not uniform: Phases 1--3 are independently machine-checked
in ProVerif~2.05 (Table~\ref{tab:proverif-results}), while Phase~4's
soundness argument is by inspection of the verification path, not a
separate ProVerif model --- a distinction the phase boxes above do not
otherwise make visible.}
\label{fig:anchor-four-phases}
\end{figure}
